\clearpage\begingroup\renewcommand{\otThumb}{}\renewcommand{\otHeadTag}{D}\phantomsection\label{dossiers}\pdfbookmark[0]{Dossiers et évaluations}{dossiers}\fiche{Une serre à construire}{Thème 1 · Modèles définis par une fonction}\phantomsection\label{projet-1}\pdfbookmark[1]{Une serre à construire}{bm-projet-1}
Une serre rectangulaire est adossée à un mur. On dispose de $40$ m de clôture pour les trois autres côtés. La largeur perpendiculaire au mur est $x$.
\repere{projet-cor-1}{Solution}
\begin{questions}
\item Exprimer l’aire et préciser le domaine.
\item Chercher les dimensions maximisant cette aire.
\item On impose une aire d’au moins $150$ m$^2$. Déterminer les largeurs possibles.
\item Présenter un schéma coté et une conclusion.
\end{questions}
\vspace{5mm}\souscours{Pour communiquer}\noindent Présenter les hypothèses, la démarche et une conclusion ; accompagner les résultats numériques de leur précision. Garder une trace de recherche, y compris des essais qui ont permis de choisir la méthode.\par\vfill\noindent{\otMedium\fontsize{8}{10}\selectfont\color{otMuted}Espace de recherche}\par\vspace{2mm}\noindent\begin{tikzpicture}\draw[otRule,line width=.2pt,step=5mm] (0,0) grid (160mm,60mm);\end{tikzpicture}
\fiche{Deux modèles de refroidissement}{Thème 2 · Modèles d’évolution}\phantomsection\label{projet-2}\pdfbookmark[1]{Deux modèles de refroidissement}{bm-projet-2}
Une température initiale vaut $80$ ; l’environnement reste à $20$. Le modèle continu est $T\,\!^{\prime}=-0{,}1(T-20)$. Le temps est en minutes.
\repere{projet-cor-2}{Solution}
\begin{questions}
\item Résoudre avec la condition initiale.
\item Construire le modèle d’Euler de pas $1$ et sa formule explicite.
\item Comparer les températures au temps $10$.
\item Quelle approximation peut-on améliorer et comment ?
\end{questions}
\vspace{5mm}\souscours{Pour communiquer}\noindent Présenter les hypothèses, la démarche et une conclusion ; accompagner les résultats numériques de leur précision. Garder une trace de recherche, y compris des essais qui ont permis de choisir la méthode.\par\vfill\noindent{\otMedium\fontsize{8}{10}\selectfont\color{otMuted}Espace de recherche}\par\vspace{2mm}\noindent\begin{tikzpicture}\draw[otRule,line width=.2pt,step=5mm] (0,0) grid (160mm,60mm);\end{tikzpicture}
\fiche{Du produit à la somme}{Thème 3 · Approche historique du logarithme}\phantomsection\label{projet-3}\pdfbookmark[1]{Du produit à la somme}{bm-projet-3}
On construit une table associant à $x$ son logarithme décimal $\ell(x)=\ln x/\ln10$.
\repere{projet-cor-3}{Solution}
\begin{questions}
\item Prouver $\ell(xy)=\ell(x)+\ell(y)$.
\item On lit $\ell(2)\simeq0{,}30103$, $\ell(3)\simeq0{,}47712$. Approcher $\ell(6)$ et $\ell(12)$.
\item Montrer que l’aire sous $1/t$ entre $1$ et $x>1$ vaut $\ln x$.
\item Expliquer les liens entre table de calcul, progression géométrique et aire sous une courbe.
\end{questions}
\vspace{5mm}\souscours{Pour communiquer}\noindent Présenter les hypothèses, la démarche et une conclusion ; accompagner les résultats numériques de leur précision. Garder une trace de recherche, y compris des essais qui ont permis de choisir la méthode.\par\vfill\noindent{\otMedium\fontsize{8}{10}\selectfont\color{otMuted}Espace de recherche}\par\vspace{2mm}\noindent\begin{tikzpicture}\draw[otRule,line width=.2pt,step=5mm] (0,0) grid (160mm,60mm);\end{tikzpicture}
\fiche{Une aire sans primitive usuelle}{Thème 4 · Calculs d’aires}\phantomsection\label{projet-4}\pdfbookmark[1]{Une aire sans primitive usuelle}{bm-projet-4}
On étudie $I=\int_0^1\mathrm e^{-x^2}\,\mathrm dx$.
\repere{projet-cor-4}{Solution}
\begin{questions}
\item Encadrer grossièrement $I$ par deux constantes.
\item Écrire des sommes de rectangles donnant un encadrement.
\item Proposer un programme Monte-Carlo dans le carré unité.
\item Comparer les garanties des deux méthodes.
\end{questions}
\vspace{5mm}\souscours{Pour communiquer}\noindent Présenter les hypothèses, la démarche et une conclusion ; accompagner les résultats numériques de leur précision. Garder une trace de recherche, y compris des essais qui ont permis de choisir la méthode.\par\vfill\noindent{\otMedium\fontsize{8}{10}\selectfont\color{otMuted}Espace de recherche}\par\vspace{2mm}\noindent\begin{tikzpicture}\draw[otRule,line width=.2pt,step=5mm] (0,0) grid (160mm,60mm);\end{tikzpicture}
\fiche{Mesurer une inégalité}{Thème 5 · Répartition des richesses}\phantomsection\label{projet-5}\pdfbookmark[1]{Mesurer une inégalité}{bm-projet-5}
On compare $L_A(x)=x^2$ et $L_B(x)=(x+x^3)/2$ pour $0\leq x\leq1$.
\repere{projet-cor-5}{Solution}
\begin{questions}
\item Vérifier que les deux modèles ont les caractéristiques d’une courbe de Lorenz.
\item Calculer les deux indices de Gini.
\item Comparer directement les courbes.
\item Une multiplication de tous les revenus par $2$ change-t-elle ces indicateurs ?
\end{questions}
\vspace{5mm}\souscours{Pour communiquer}\noindent Présenter les hypothèses, la démarche et une conclusion ; accompagner les résultats numériques de leur précision. Garder une trace de recherche, y compris des essais qui ont permis de choisir la méthode.\par\vfill\noindent{\otMedium\fontsize{8}{10}\selectfont\color{otMuted}Espace de recherche}\par\vspace{2mm}\noindent\begin{tikzpicture}\draw[otRule,line width=.2pt,step=5mm] (0,0) grid (160mm,60mm);\end{tikzpicture}
\fiche{Une information qui change la décision}{Thème 6 · Inférence bayésienne}\phantomsection\label{projet-6}\pdfbookmark[1]{Une information qui change la décision}{bm-projet-6}
Deux fournisseurs $A,B$ livrent respectivement $80\,\%$ et $20\,\%$ des composants. Les défauts concernent $1\,\%$ des composants de $A$ et $6\,\%$ de ceux de $B$.
\repere{projet-cor-6}{Solution}
\begin{questions}
\item Construire un arbre et calculer le taux de défaut.
\item Quelle est la probabilité qu’un composant défectueux vienne de $B$ ?
\item À quelle part de marché $p$ de $B$ les défauts provenant de $B$ deviennent-ils majoritaires parmi tous les défauts ?
\item Expliquer la différence entre probabilité a priori et a posteriori.
\end{questions}
\vspace{5mm}\souscours{Pour communiquer}\noindent Présenter les hypothèses, la démarche et une conclusion ; accompagner les résultats numériques de leur précision. Garder une trace de recherche, y compris des essais qui ont permis de choisir la méthode.\par\vfill\noindent{\otMedium\fontsize{8}{10}\selectfont\color{otMuted}Espace de recherche}\par\vspace{2mm}\noindent\begin{tikzpicture}\draw[otRule,line width=.2pt,step=5mm] (0,0) grid (160mm,60mm);\end{tikzpicture}
\fiche{Contrôler une surréservation}{Thème 7 · Répétitions et échantillonnage}\phantomsection\label{projet-7}\pdfbookmark[1]{Contrôler une surréservation}{bm-projet-7}
Une petite navette accepte $12$ réservations pour $10$ places. Chaque client se présente indépendamment avec probabilité $0{,}8$.
\repere{projet-cor-7}{Solution}
\begin{questions}
\item Identifier la loi du nombre $X$ de présents.
\item Calculer $P(X>10)$.
\item La règle « risque de dépassement au plus $0{,}05$ » est-elle respectée ?
\item Citer deux hypothèses dont la violation modifierait la conclusion.
\end{questions}
\vspace{5mm}\souscours{Pour communiquer}\noindent Présenter les hypothèses, la démarche et une conclusion ; accompagner les résultats numériques de leur précision. Garder une trace de recherche, y compris des essais qui ont permis de choisir la méthode.\par\vfill\noindent{\otMedium\fontsize{8}{10}\selectfont\color{otMuted}Espace de recherche}\par\vspace{2mm}\noindent\begin{tikzpicture}\draw[otRule,line width=.2pt,step=5mm] (0,0) grid (160mm,60mm);\end{tikzpicture}
\fiche{Attendre sans vieillir ?}{Thème 8 · Temps d’attente}\phantomsection\label{projet-8}\pdfbookmark[1]{Attendre sans vieillir ?}{bm-projet-8}
Un composant fictif a une durée $T$ exponentielle de moyenne $50$ heures. On le compare à une attente $U$ uniforme sur $[0,100]$.
\repere{projet-cor-8}{Solution}
\begin{questions}
\item Donner les deux espérances.
\item Calculer $P(T>40\mid T>10)$.
\item Calculer $P(U>40\mid U>10)$.
\item Comparer avec les probabilités d’attendre encore plus de $30$ heures au départ.
\end{questions}
\vspace{5mm}\souscours{Pour communiquer}\noindent Présenter les hypothèses, la démarche et une conclusion ; accompagner les résultats numériques de leur précision. Garder une trace de recherche, y compris des essais qui ont permis de choisir la méthode.\par\vfill\noindent{\otMedium\fontsize{8}{10}\selectfont\color{otMuted}Espace de recherche}\par\vspace{2mm}\noindent\begin{tikzpicture}\draw[otRule,line width=.2pt,step=5mm] (0,0) grid (160mm,60mm);\end{tikzpicture}
\fiche{Une croissance observée}{Thème 9 · Corrélation et causalité}\phantomsection\label{projet-9}\pdfbookmark[1]{Une croissance observée}{bm-projet-9}
Les couples fictifs $(t,y)$ sont $(0,2),(1,4),(2,8),(3,16)$.
\repere{projet-cor-9}{Solution}
\begin{questions}
\item Construire le nuage et proposer un modèle.
\item Linéariser par logarithme et retrouver le modèle.
\item Prévoir $y$ pour $t=4$ puis $t=20$ ; comparer la portée des prévisions.
\item Cette association exacte prouve-t-elle que le temps est une cause du phénomène ?
\end{questions}
\vspace{5mm}\souscours{Pour communiquer}\noindent Présenter les hypothèses, la démarche et une conclusion ; accompagner les résultats numériques de leur précision. Garder une trace de recherche, y compris des essais qui ont permis de choisir la méthode.\par\vfill\noindent{\otMedium\fontsize{8}{10}\selectfont\color{otMuted}Espace de recherche}\par\vspace{2mm}\noindent\begin{tikzpicture}\draw[otRule,line width=.2pt,step=5mm] (0,0) grid (160mm,60mm);\end{tikzpicture}
\fiche{Évaluation 1 · Analyse}{60 minutes · 20 points}\phantomsection\label{evaluation-1}\pdfbookmark[1]{Évaluation 1 · Analyse}{bm-evaluation-1}
Pour $t\geq0$, $f(t)=20+80\mathrm e^{-0{,}25t}$.
\repere{evaluation-cor-1}{Solution}
\begin{questions}
\item Calculer $f(0)$, la dérivée et la limite en $+\infty$. \hfill(4 points)
\item Résoudre $f(t)\leq40$. \hfill(4 points)
\item Donner l’équation différentielle et la condition initiale de ce modèle. \hfill(4 points)
\item Calculer la moyenne de $f$ sur $[0,4]$. \hfill(4 points)
\item Construire Euler avec pas $1$ ; calculer deux étapes et comparer à l’exact. \hfill(4 points)
\end{questions}
\vspace{5mm}\souscours{Pour communiquer}\noindent Présenter les hypothèses, la démarche et une conclusion ; accompagner les résultats numériques de leur précision. Garder une trace de recherche, y compris des essais qui ont permis de choisir la méthode.\par\vfill\noindent{\otMedium\fontsize{8}{10}\selectfont\color{otMuted}Espace de recherche}\par\vspace{2mm}\noindent\begin{tikzpicture}\draw[otRule,line width=.2pt,step=5mm] (0,0) grid (160mm,60mm);\end{tikzpicture}
\fiche{Évaluation 2 · Probabilités}{60 minutes · 20 points}\phantomsection\label{evaluation-2}\pdfbookmark[1]{Évaluation 2 · Probabilités}{bm-evaluation-2}
Un atelier $A$ fabrique $70\,\%$ des pièces, dont $2\,\%$ sont défectueuses. L’atelier $B$ fabrique le reste, dont $5\,\%$ sont défectueuses.
\repere{evaluation-cor-2}{Solution}
\begin{questions}
\item Calculer la probabilité $d$ d’un défaut. \hfill(4 points)
\item Une pièce est défectueuse : calculer la probabilité qu’elle vienne de $B$. \hfill(4 points)
\item Pour dix pièces indépendantes au même taux $d$, identifier la loi du nombre de défauts et la probabilité d’au moins un défaut. \hfill(4 points)
\item On répète des contrôles indépendants jusqu’au premier défaut. Donner la loi du rang, l’espérance et la probabilité de plus de vingt contrôles. \hfill(4 points)
\item Une attente continue suit une exponentielle de moyenne $30$ minutes. Calculer $P(T>15)$ et $P(T>25\mid T>10)$. \hfill(4 points)
\end{questions}
\vspace{5mm}\souscours{Pour communiquer}\noindent Présenter les hypothèses, la démarche et une conclusion ; accompagner les résultats numériques de leur précision. Garder une trace de recherche, y compris des essais qui ont permis de choisir la méthode.\par\vfill\noindent{\otMedium\fontsize{8}{10}\selectfont\color{otMuted}Espace de recherche}\par\vspace{2mm}\noindent\begin{tikzpicture}\draw[otRule,line width=.2pt,step=5mm] (0,0) grid (160mm,60mm);\end{tikzpicture}
\fiche{Évaluation 3 · Statistique et inégalités}{60 minutes · 20 points}\phantomsection\label{evaluation-3}\pdfbookmark[1]{Évaluation 3 · Statistique et inégalités}{bm-evaluation-3}
Les observations fictives sont $(1,2),(2,3),(3,5),(4,6)$. On considère aussi la répartition $L(x)=x^2$.
\repere{evaluation-cor-3}{Solution}
\begin{questions}
\item Calculer le point moyen et les quantités nécessaires à une régression affine. \hfill(4 points)
\item Déterminer la droite et prédire en $x=3{,}5$, puis en $x=10$. Discuter. \hfill(4 points)
\item Calculer le coefficient de corrélation ; commenter sans conclure à une causalité. \hfill(4 points)
\item Vérifier que $L$ est une courbe de Lorenz et interpréter $L(0{,}6)$. \hfill(4 points)
\item Calculer Gini et expliquer l’effet d’une multiplication de tous les revenus par une même constante positive. \hfill(4 points)
\end{questions}
\vspace{5mm}\souscours{Pour communiquer}\noindent Présenter les hypothèses, la démarche et une conclusion ; accompagner les résultats numériques de leur précision. Garder une trace de recherche, y compris des essais qui ont permis de choisir la méthode.\par\vfill\noindent{\otMedium\fontsize{8}{10}\selectfont\color{otMuted}Espace de recherche}\par\vspace{2mm}\noindent\begin{tikzpicture}\draw[otRule,line width=.2pt,step=5mm] (0,0) grid (160mm,60mm);\end{tikzpicture}
\clearpage\endgroup